
\documentclass{article}
\usepackage{colortbl}
\usepackage{makecell}
\usepackage{multirow}
\usepackage{supertabular}

\begin{document}

\newcounter{utterance}

\centering \large Interaction Transcript for game `eqbench', experiment `eq\_bench', episode 31 with reflexion{-}Qwen3.5{-}9B.
\vspace{24pt}

{ \footnotesize  \setcounter{utterance}{1}
\setlength{\tabcolsep}{0pt}
\begin{supertabular}{c@{$\;$}|p{.15\linewidth}@{}p{.15\linewidth}p{.15\linewidth}p{.15\linewidth}p{.15\linewidth}p{.15\linewidth}}
    \# & \multicolumn{2}{c}{Player} && \multicolumn{2}{c}{Game Master} \\
    \hline

    \theutterance \stepcounter{utterance}  
    & & & \multicolumn{4}{p{0.6\linewidth}}{
        \cellcolor[rgb]{0.9,0.9,0.9}{
            \makecell[{{p{\linewidth}}}]{
                \texttt{\tiny{[P1$\langle$GM]}}
                \\ 
\texttt{Your task is to predict the likely emotional responses of a character in this dialogue:} \\
\\ 
\texttt{Brenda: So, you thought it was a good idea to turn our home into a nightclub?} \\
\texttt{Tessa: It wasn't a nightclub, Mom. Just a few friends hanging out.} \\
\texttt{Brenda: Hanging out? The neighbors called the police, Tessa!} \\
\texttt{Tessa: Maybe if you weren't so uptight, they wouldn't have to.} \\
\texttt{Brenda: This is not about me being uptight. This is about you respecting our home.} \\
\texttt{Tessa: Our home? Or your home?} \\
\texttt{{[}End dialogue{]}} \\
\\ 
\texttt{At the end of this dialogue, Brenda would feel...} \\
\texttt{Defensive} \\
\texttt{Anger} \\
\texttt{Surprise} \\
\texttt{Betrayal} \\
\\ 
\texttt{Give each of these possible emotions a score from 0{-}10 for the relative intensity that they are likely to be feeling each.} \\
\\ 
\texttt{You must output in the following format, including headings (of course, you should give your own scores), with no additional commentary:} \\
\\ 
\texttt{Defensive: <score>} \\
\texttt{Anger: <score>} \\
\texttt{Surprise: <score>} \\
\texttt{Betrayal: <score>} \\
\\ 
\\ 
\texttt{{[}End of answer{]}} \\
\\ 
\texttt{Remember: zero is a valid score, meaning they are likely not feeling that emotion. You must score at least one emotion > 0.} \\
\\ 
\texttt{Your answer:} \\
            }
        }
    }
     \\ \\

    \theutterance \stepcounter{utterance}  
    & \multicolumn{4}{p{0.6\linewidth}}{
        \cellcolor[rgb]{0.561,0.737,0.561}{
            \makecell[{{p{\linewidth}}}]{
                \texttt{\tiny{[P1$\rangle$GM]}}
                \texttt{Defensive: 8} \\
\texttt{Anger: 7} \\
\texttt{Surprise: 2} \\
\texttt{Betrayal: 3} \\
            }
        }
    }
    & & \\ \\

    \theutterance \stepcounter{utterance}  
    & & & \multicolumn{2}{p{0.3\linewidth}}{
        \cellcolor[rgb]{0.9,0.9,0.9}{
            \makecell[{{p{\linewidth}}}]{
                \texttt{\tiny{[GM$|$GM]}}
                \texttt{Defensive: 8} \\
\texttt{Anger: 7} \\
\texttt{Surprise: 2} \\
\texttt{Betrayal: 3} \\
            }
        }
    }
    & & \\ \\

    \theutterance \stepcounter{utterance}  
    & & & \multicolumn{2}{p{0.3\linewidth}}{
        \cellcolor[rgb]{0.9,0.9,0.9}{
            \makecell[{{p{\linewidth}}}]{
                \texttt{\tiny{[GM$|$GM]}}
                \texttt{Defensive: 7} \\
\texttt{Anger: 9} \\
\texttt{Surprise: 0} \\
\texttt{Betrayal: 6} \\
            }
        }
    }
    & & \\ \\

    \theutterance \stepcounter{utterance}  
    & & & \multicolumn{2}{p{0.3\linewidth}}{
        \cellcolor[rgb]{0.9,0.9,0.9}{
            \makecell[{{p{\linewidth}}}]{
                \texttt{\tiny{[GM$|$GM]}}
                \texttt{game\_result = LOSE} \\
            }
        }
    }
    & & \\ \\

    \theutterance \stepcounter{utterance}  
    & & \multicolumn{4}{p{0.6\linewidth}}{
        \cellcolor[rgb]{0.95,0.95,0.95}{
            \makecell[{{p{\linewidth}}}]{
                \texttt{\tiny{[P1$\leftrightarrow$P]}}
                \texttt{Post{-}episode Reflexion (not sent to GM):} \\
\texttt{New plan: I failed because I assigned high scores to multiple emotions simultaneously (Defensive, Anger, Surprise, Betrayal) without prioritizing the single most dominant response based on the dialogue's climax. Brenda's primary reaction to Tessa's accusation ("Or your home?") is a defensive reassertion of authority and ownership, not a shift to anger or surprise. Next time, I will analyze the final exchange to identify the core emotional driver, assign a high score (8{-}10) to that specific emotion, and assign zero or very low scores to the others to ensure a clear, singular prediction. Specifically, I will recognize that Tessa's challenge to Brenda's authority triggers a defensive posture, making "Defensive" the only high{-}score option.} \\
            }
        }
    }
    & \\ \\

\end{supertabular}
}

\end{document}
